\chapter[Introduction]{Introduction}
\label{Chap:Intro}

% ***************************************************
% Introduction
% ***************************************************

\noindent
Static application security testing (SAST) is the application of static program analysis to the security of software projects. MORE INFO ABOUT STATIC ANALYSIS HERE. Static analysis is performed on the source code, compiled byte-code, or some other representation of the source code \cite{NielsonFlemming2005Popa}. Popular SAST tools include: SonarQube \cite{sonarqube}, SpotBugs \cite{spotbugs}, Snyk Code \cite{snykcode}, and Facebook's Infer \cite{infer}. SAST tools have become widely used at all stages of the software development life-cycle. Furthermore, they have been demonstrated to be highly effective in identifying defects in code compared to other methods, such as manual inspection \cite{Nichols2020}.\\

\noindent
Unfortunately, SAST tools are not perfect. From Rice's Theorem, it is known that any non-trivial property about programs is reducible to the halting problem. That is, it is undecidable --- a "perfect" analysis is not always possible \cite{moller2012static}. As such, SAST tools are required to make approximations, resulting in an imperfect output --- i.e. false positives and false negatives \cite{chess2004static}. False positives are not ideal as manual inspection is required to separate genuine defects from false positives, leading to increased time and labour cost. And, to reduce the rate of false negatives, manual configuration of the SAST tools for specific software projects is often required \cite{ZampettiFiorella2017HOSP}.\\

\noindent
For example, to perform an effective taint analysis a custom configuration of sources, sinks and sanitisers for a particular project is often required when the default configuration proves insufficient. Creating this manual configuration is a time consuming and resource intensive process. As such, methodologies to automatically label methods as sources, sinks, and sanitisers may enable existing taint analysis tools to operate more effectively and efficiently. It should also be noted that in cases where a closed-source library is being used in a software project, it may be impossible to determine which methods in the API should be manually added to the taint analysis.\\

\noindent
Deep learning, a subtopic within the broader machine learning, has seen an astounding rise in popularity in recent years. What was once a niche topic reserved for academics with access to computational resources far out of reach by most has turned into arguably one of the most influential technologies of the modern age. Deep learning has fundamentally changed our modern life, with applications ranging from targeted advertising to identifying COVID-19 cases from lung scans \cite{RAHIMZADEH2021102588}. One often over-looked application of deep-learning, lost amidst hype concerning self-driving cars and ... is natural-language processing (NLP).\\

\noindent
NLP has existed in one form or another for decades. In the early days of NLP, work was painstakingly done using handwritten rules before moving into probabilistic methods in the 1990s. Naturally, NLP was completely revolutionised with the rising popularity of deep-learning \cite{DengLi2018Dlin}. Furthermore, with companies such as Google and Microsoft publicly releasing large pre-trained models, NLP has never been more accessible. Popular applications of NLP (using deep-learning) is machine translation and speech recognition. One application of deep-learning, and NLP, that has not been explored thoroughly in the literature is its application to software security.\\

\noindent
Javadoc, while being an informal description of code, may contain information relevant to software security. With powerful pre-trained models such as CodeBERT, which \blockquote{captures the semantic connection between natural language and programming language, and produces general-purpose
representations that can broadly support NL-PL understanding tasks}\cite{CodeBERT}, it may be possible to fine-tune CodeBERT to extract security relevant properties from Javadocs in order to configure existing SAST tools.\\

\section{Topic Definition}
%Replace "Your thesis topic" with what will become the \label{} for this section.

% ********* Enter your text below this line: ********
Introduce your topic.

% ***************************************************